% Propietario conservado: /Users/ruben/Documents/New project/output/EXTREMA_MEDIA_RAZON_RECIPROCIDAD_ES_EN_20260918/ARTICULO/ES/source/nuclear/11.tex
\chapter{Reconstrucción dimensional y conservación exterior}
\label{x_reciprocidad:nuclear:11}
\section{Dos memorias recuperan el registro centrado}

\subsection{Operadores de lectura del registro}

En la carta ordenada del registro \(K\) de
\cref{x_reciprocidad:sec:k-registro,x_reciprocidad:img09:imagen}, sea \(V=\RR^{12}\) su
realización lineal posterior y sea \(S\) la permutación
\((Sx)_i=x_{i+1}\), con índices módulo \(12\).
Los lectores \(I-S^3\) e \(I-S^4\) son los del extractor transversal
de \cref{x_reciprocidad:img09:imagen}.
La compresión nonádica de \cref{x_reciprocidad:nuclear:10} proporciona,
para \(j=3,4\),



\begin{equation}
T_j=\frac{8I+S^j}{9},\qquad D_j=\frac{\sqrt8}{9}(S^j-I),
\qquad T_j^*T_j+D_j^*D_j=I.
\label{x_reciprocidad:nuc11:eq:1}
\end{equation}



La partición \(8+1\) y la normalización \(9\) proceden de las nueve
posiciones y de su contraste colectivo.
Se considera la siguiente composición explícita de lectores,
aplicada al registro generado:



\begin{equation}
m_0=D_3K,\qquad m_1=D_4T_3K,
\qquad M K=(m_0,m_1).
\label{x_reciprocidad:nuc11:eq:2}
\end{equation}



La segunda memoria actúa sobre la lectura producida por la primera
etapa. Los complementos se archivan por separado.
La evolución completa con reentrada corresponde al acoplamiento
reversible de \cref{x_reciprocidad:nuclear:10}.

\subsection{Inversa exacta}

Puesto que \((S^3)^4=I\),



\begin{equation}
T_3^{-1}=\frac{512I-64S^3+8S^6-S^9}{455}.
\label{x_reciprocidad:nuc11:eq:3}
\end{equation}



En efecto,
\((8I+S^3)(512I-64S^3+8S^6-S^9)=4095I\)
y \(4095=9\cdot455\).
Por tanto, las memorias proporcionan



\begin{equation}
b=-\frac9{\sqrt8}m_0=(I-S^3)K,
\quad c=-\frac9{\sqrt8}T_3^{-1}m_1=(I-S^4)K.
\label{x_reciprocidad:nuc11:eq:4}
\end{equation}



La conmutación de \(T_3\) con \(S^4\) justifica la segunda igualdad.
Definamos



\[
w_i=c_i-b_{i+1},\quad B_0=0,\quad
B_i=\sum_{j<i}w_j.
\]



La identidad \(w_i=K_i-K_{i+1}\) prueba que \(B_i=K_0-K_i\).
Si \(q=\sum_{i=0}^{11}K_i\), se recupera exactamente



\begin{equation}
K_i=\frac{q+\sum_{j=0}^{11}B_j}{12}-B_i.
\label{x_reciprocidad:nuc11:eq:5}
\end{equation}



Es la inversa \eqref{x_reciprocidad:img09:inversa}, compuesta ahora con la
compresión y la memoria. Con \(q=0\), la fórmula devuelve
\(P_{11}K\), donde



\[
P_{11}=I-\frac1{12}\mathbf1\mathbf1^*.
\]



Las dos memorias recuperan las once componentes de media nula.
La carga uniforme se recupera mediante \(q\).
El núcleo de \(M\) es precisamente la recta uniforme:
anular ambas diferencias implica fijación por \(S^3\) y \(S^4\);
como \(S=S^4(S^3)^{-1}\), implica fijación por \(S\).

\subsection{Cota óptima de estabilidad}

El Gram del procedimiento es



\begin{equation}
G=M^*M=D_3^*D_3+T_3^*D_4^*D_4T_3
=I-(T_4T_3)^*(T_4T_3).
\label{x_reciprocidad:nuc11:eq:6}
\end{equation}



En el modo \(k\) de \(S\), el valor propio de \(T_j^*T_j\) es
\[
 \frac{65+16\cos(2\pi jk/12)}{81}.
\]
Su sustitución en el Gram anterior da, con multiplicidades,



\begin{equation}
\operatorname{Spec}G=\left\{
0^{[1]},\ (16/81)^{[2]},\ (8/27)^{[2]},\ (32/81)^{[1]},
(952/2187)^{[4]},\ (1256/2187)^{[2]}\right\}.
\label{x_reciprocidad:nuc11:eq:7}
\end{equation}



El modo nulo es la media. Sea \(L\) la inversa de mínimos cuadrados
de \(M\) sobre \(\mathbf1^\perp\), extendida al espacio de
registros por



\[
L=(G|_{\mathbf1^\perp})^{-1}M^*.
\]



Entonces



\begin{equation}
LM=P_{11},\qquad \|L\|=\frac94.
\label{x_reciprocidad:nuc11:eq:8}
\end{equation}



La cota es óptima porque los modos \(k=3,9\) alcanzan \(16/81\).
Para datos perturbados y reconstrucción con \(L\),



\begin{equation}
\|\delta K\|^2\le \frac{|\delta q|^2}{12}
+\frac{81}{16}\bigl(\|\delta m_0\|^2+\|\delta m_1\|^2\bigr).
\label{x_reciprocidad:nuc11:eq:9}
\end{equation}



Las componentes uniforme y centrada son ortogonales. La fórmula \eqref{x_reciprocidad:nuc11:eq:5} constituye una inversa exacta sobre datos compatibles; \eqref{x_reciprocidad:nuc11:eq:8} especifica la extensión estable usada para datos con ruido, que no tienen por qué satisfacer las compatibilidades de \eqref{x_reciprocidad:nuc11:eq:5}.

\subsection{Ejemplo con el registro dodecafásico generado}

El registro de \eqref{x_reciprocidad:eq:k-vector} es



\begin{equation}
K=(234,543,140,729,659,824,621,58,914,794,146,601).
\label{x_reciprocidad:nuc11:eq:10}
\end{equation}



Es una salida del generador anterior, utilizada como caso de
aplicación. Los coeficientes y los lectores se fijan antes de
evaluar este registro. Para almacenar las memorias con aritmética
racional, escribimos \(z_i=m_i/\sqrt8\).
Los dos registros efectivos son



\[
9z_0=(495,116,684,-108,-601,90,173,88,-313,-560,397,-461),
\]





\begin{equation}
81z_1=(2729,2503,3818,-5843,2583,-920,-4051,4338,-5312,-1583,233,1505).
\label{x_reciprocidad:nuc11:eq:11}
\end{equation}



La sustitución de estos registros en las fórmulas de reconstrucción,
con \(q=6263\), reproduce las doce componentes de \(K\).
Con \(q=0\) reproduce \(K-(6263/12)\mathbf1\).
El certificado suplementario reconstruye primero el registro desde
las memorias y después lo compara con \eqref{x_reciprocidad:eq:k-vector}.

Añadir una tercera memoria \(D_3T_4T_3K\) eleva el menor valor
propio positivo del Gram a \(8/27\).
La amplificación óptima desciende de \(9/4\) a \(\sqrt{27/8}\).
Una tercera memoria \(D_4T_4T_3K\) mantiene \(16/81\).
La elección entre repetir el paso \(3\) o el paso \(4\) tiene
un efecto de estabilidad calculable antes de evaluar el registro.

\section{Reconstrucción del proyector dimensional desde dos memorias}

\subsection{Composición con la incidencia ya construida}

La carta excepcional de \cref{x_reciprocidad:exc:k-direccion} define la
representación ortogonal de \(A_5\) sobre las doce clases de
\(A_5/C_5\), con representantes y carácter explícitos.
Su proyector isótipo tridimensional es
\eqref{x_reciprocidad:eq:k-proyectores-incidencia}:



\begin{equation}
P_3=\frac3{60}\sum_{a\in A_5}\chi_3(a^{-1})\rho(a),
\quad P_3P_{11}=P_3,\quad P_3\mathbf1=0.
\label{x_reciprocidad:nuc11:eq:12}
\end{equation}



El registro generado precede a esta representación.
El \cref{x_reciprocidad:lem:k-sector-no-nulo} y
\eqref{x_reciprocidad:eq:k-norma-direccion} y \eqref{x_reciprocidad:eq:k-proyector-diez} establecen



\[
u_K=P_3P_{11}K,\qquad
\|u_K\|^2=\frac{6638585+2275584\sqrt5}{20}>0,
\]





\begin{equation}
P_{10}(K)=P_{11}-\frac{u_Ku_K^*}{\|u_K\|^2}.
\label{x_reciprocidad:nuc11:eq:13}
\end{equation}



La composición nueva de \eqref{x_reciprocidad:nuc11:eq:8}, \eqref{x_reciprocidad:nuc11:eq:12} y \eqref{x_reciprocidad:nuc11:eq:13} es



\begin{equation}
\boxed{
u_K=P_3L(m_0,m_1),\qquad
P_{10}(K)=P_{11}-
\frac{P_3L(m)\,[P_3L(m)]^*}{\|P_3L(m)\|^2}.
}
\label{x_reciprocidad:nuc11:eq:14}
\end{equation}



\begin{proof}
Las identidades \(LM=P_{11}\) y \(P_3P_{11}=P_3\) dan
\(P_3LMK=P_3K=u_K\).
La no nulidad está evaluada en \eqref{x_reciprocidad:eq:k-norma-direccion}.
El cociente de rango uno coincide, por tanto, con el de
\eqref{x_reciprocidad:eq:k-proyector-diez}.
La composición recupera el proyector ortogonal íntegro de rango
diez a partir de las dos memorias, con independencia de la carga
\(q\) y de la lectura terminal \(T_4T_3K\).
\end{proof}

La reducción \(12\to11\) elimina la dirección uniforme y la
reducción \(11\to10\) utiliza el vector seleccionado por \(K\).
La composición recupera la matriz completa del proyector en el
portador lineal de incidencia.
Sus realizaciones dimensionales mantienen los mapas de
\cref{x_reciprocidad:prop:k-reduccion-diez}.

\subsection{Error del proyector}

Sea \(\varepsilon\) la norma conjunta del error de las dos
memorias y sea \(\widehat u=P_3L(m+\delta m)\).
Entonces \(\|\widehat u-u_K\|\le9\varepsilon/4\).
Si \(9\varepsilon/4<\|u_K\|\), el vector \(\widehat u\) es no
nulo y



\begin{equation}
\boxed{
\|\widehat P_{10}-P_{10}\|
\le\frac{9\varepsilon}{4\|u_K\|}.
}
\label{x_reciprocidad:nuc11:eq:15}
\end{equation}



\begin{proof}
La distancia en norma operatoria entre dos proyectores ortogonales
sobre rectas es el seno del ángulo entre ellas.
La distancia de \(u_K\) a la recta de \(\widehat u\) es a lo sumo
\(\|u_K-\widehat u\|\). Dividir por \(\|u_K\|\) prueba la cota.
Aquí \(\|u_K\|\approx765{,}733\), y el coeficiente es
aproximadamente \(0{,}00293836\).
\end{proof}
Para registros próximos a \(\ker P_3\), la estabilidad de la
dirección depende de \(\|u_K\|\).
La cota \(9/4\) para el registro centrado permanece válida.

\subsection{Un descriptor escalar de K, con invariancias explícitas}

Se distinguen el registro \(K\), sus lecturas escalares y sus
invariantes. Para \(K\notin\RR\mathbf1\), definimos



\begin{equation}
\eta_K=\frac{\|P_3P_{11}K\|^2}{\|P_{11}K\|^2}.
\label{x_reciprocidad:nuc11:eq:16}
\end{equation}



Para el registro \eqref{x_reciprocidad:eq:k-vector},
\(\sum_iK_i^2=4251257\) y
\(\|P_{11}K\|^2=11789915/12\). Así,



\begin{equation}
\eta_K=
\frac{3(6638585+2275584\sqrt5)}{5\cdot11789915}
=\frac{13089003+13653504\varphi_{\rm HMT}}{58949575}
\approx0,5967954228259091.
\label{x_reciprocidad:nuc11:eq:17}
\end{equation}



La segunda escritura utiliza el reconocimiento cuadrático posterior
de \(\varphi_{\rm HMT}\).
Se cumple \(0\le\eta_K\le1\).
El cociente es invariante bajo
\(K\mapsto aK+b\mathbf1\), con \(a\ne0\), y bajo cambios
ortogonales simultáneos de \(K,P_{11},P_3\).
En la carta fija es invariante bajo \(\rho(A_5)\), pues ambos
proyectores conmutan con esa representación.
Es recuperable desde las memorias porque ambas normas que lo
definen lo son.
Mide la fracción cuadrática del registro centrado situada en el
sector isótipo elegido.
El registro completo conserva además la información que este
descriptor escalar identifica.
La procedencia de esta formalización adicional se documenta en
el suplemento; las invariancias afines aquí probadas pertenecen
a la realización lineal.

\section{Aplicación sobre un vector efectivo de la construcción excepcional}

La construcción de
\cref{x_reciprocidad:km:orden-tres-reticular,x_reciprocidad:km:estabilidad-vecino}
utiliza doce fibras \(A_2\), su pegado y el vecino reticular.
La palabra trítica es
\(c^\star=(1,1,1,1,1,1,2,1,1,1,1,1)\).
Sus bloques actúan por \(C^{-1}\) cuando la componente es \(1\)
y por \(C\) cuando es \(2\), en las cartas especificadas por
\cref{x_reciprocidad:km:palabra-total,x_reciprocidad:km:isometria-ternaria}.
El operador \(g\) resultante cumple
\(g^3=I\), \(g^2+g+I=0\), y es ortogonal para la suma de las
métricas \(A_2\).
El descenso al pegado y al vecino se demuestra en
\cref{x_reciprocidad:km:estabilidad-vecino}.
Los cálculos siguientes actúan sobre su portador real de
dimensión \(24\).

En cada fibra,
\[
 C=\begin{pmatrix}0&-1\\1&-1\end{pmatrix},
 \qquad
 G_{A_2}=\begin{pmatrix}2&-1\\-1&2\end{pmatrix}.
\]
Sobre la suma de las doce fibras se define



\[
T=\frac{8I+g}{9},\qquad D=\frac{\sqrt8}{9}(g-I).
\]



Con el adjunto relativo a esa métrica,
\(g^*=g^{-1}=-g-I\), y se obtiene



\begin{equation}
T^*T=\frac{19}{27}I,
\qquad D^*D=\frac8{27}I,
\qquad D^{-1}=-\frac3{\sqrt8}(g+2I).
\label{x_reciprocidad:nuc11:eq:18}
\end{equation}



La inversa procede de \((g-I)(g+2I)=-3I\).
La construcción del vecino fija el vector de
\eqref{x_reciprocidad:exc:vector},
\(v=4\rho_1+\rho_2+\cdots+\rho_{12}\), con una raíz
\(\rho_b\) de norma cuadrática \(2\) en cada fibra.
Por tanto,



\begin{equation}
\|v\|^2=54,\qquad \|Tv\|^2=38,\qquad \|Dv\|^2=16,
\qquad v=-\frac3{\sqrt8}(g+2I)Dv.
\label{x_reciprocidad:nuc11:eq:19}
\end{equation}



Un solo complemento recupera este vector completo.
La lectura conserva \(38\) unidades de norma cuadrática y el
registro contiene \(16\); su suma es \(54\).
El resultado utiliza el vector de \eqref{x_reciprocidad:exc:vector} y sus fibras
reticulares.

\section{Ley de conservación en cada grado exterior}

\subsection{Enunciado y prueba}

Sean \(P_0=I\), \(P_n=T_{n-1}\cdots T_0\) y



\[
\mathcal J_Nv=(P_Nv,D_0P_0v,\ldots,D_{N-1}P_{N-1}v).
\]



La ley de \cref{x_reciprocidad:nuclear:8,x_reciprocidad:nuclear:10} conserva la forma \(G_0\)
mediante la suma ortogonal de las formas correspondientes:
\(\mathcal J_N^*G_{\rm sal}\mathcal J_N=G_0\).
Para cada grado exterior finito admisible \(p\), se deduce



\begin{equation}
\boxed{
(\Lambda^p\mathcal J_N)^*(\Lambda^pG_{\rm sal})
(\Lambda^p\mathcal J_N)=\Lambda^pG_0.
}
\label{x_reciprocidad:nuc11:eq:20}
\end{equation}



La forma inducida se define por



\[
\langle x_1\wedge\cdots\wedge x_p,
y_1\wedge\cdots\wedge y_p\rangle_{\Lambda^pG}
=\det(\langle x_i,y_j\rangle_G)_{i,j=1}^p.
\]



\begin{proof}
Sustituir el balance de grado uno en cada entrada de ese determinante. La igualdad se cumple sobre vectores descomponibles y se extiende bilinealmente porque éstos generan la potencia exterior.
\end{proof}

La descomposición relevante es



\[
\Lambda^p\Bigl(\bigoplus_aH_a\Bigr)
\simeq\bigoplus_{\sum n_a=p}\bigotimes_a\Lambda^{n_a}H_a.
\]



En el caso positivo, si
\(\mathcal J_{p,\mathbf n}\) designa las componentes de ese mapa,



\begin{equation}
\boxed{\|\xi\|^2=\sum_{\sum n_a=p}
\|\mathcal J_{p,\mathbf n}\xi\|^2.}
\label{x_reciprocidad:nuc11:eq:21}
\end{equation}



Se conservan las componentes puramente visibles, las puramente
registradas y las mixtas.
Para \(p=2\) y \(p=3\), se conservan, respectivamente, las normas
cuadráticas de área y volumen.
Para una forma alternante o indefinida, la identidad exterior
conserva su interpretación bilineal.

\subsection{Distribución exacta entre los sectores exteriores}

En la realización Coxeter anterior,
\(a=19/27\) y \(b=8/27\).
Normalizar \(T\) y \(D\) da dos isometrías, de modo que la
componente con \(k\) direcciones en la memoria tiene norma cuadrática



\begin{equation}
\boxed{
\|\mathcal J_{p,k}\xi\|^2=
\binom pk a^{p-k}b^k\|\xi\|^2.
}
\label{x_reciprocidad:nuc11:eq:22}
\end{equation}



Se prueba expandiendo exteriormente
\(x\mapsto(\sqrt a\,x,\sqrt b\,x)\).
Las componentes con posiciones distintas de memoria son
ortogonales; las isometrías normalizadas \(T\) y \(D\)
conservan sus normas.

Para un bivector de norma uno,

\begin{center}
\begin{tabular}{@{}lr@{}}
\toprule
Sector & Norma cuadrática\\
\midrule
Dos direcciones en la lectura & \(361/729\)\\
Una dirección en lectura y otra en memoria & \(304/729\)\\
Dos direcciones en memoria & \(64/729\)\\
\bottomrule
\end{tabular}
\end{center}

La suma es uno. El sector mixto contiene exactamente \(304/729\),
aproximadamente \(41{,}7\%\), de la identidad de área al cuadrado.
En grado tres, las cuatro componentes son
\((6859,8664,3648,512)/19683\); las componentes mixtas suman
\(12312/19683\).

La ley se aplica a \(0\le p\le24\) en el portador de las doce
fibras \(A_2\), y functorialmente a otros portadores generados.
El grado exterior, la resolución de refinamiento, la escala de
una métrica y la dimensión física poseen tipados distintos.
El cambio \(G\mapsto s^2G\), con \(s>0\), multiplica la norma
cuadrática de grado \(p\) por \(s^{2p}\); ambas partes de la
identidad exterior se transforman de igual manera.
La realización de las dimensiones \(10\) y \(11\) conserva los
mapas dimensionales previamente especificados.

\subsection{Naturalidad y profundidad arbitraria}

El refinamiento de historias generado conserva la medida mediante



\[
R\delta_h=\sum_{\epsilon\in\operatorname{Out}(h)}
\sqrt{p_h(\epsilon)}\,\delta_{h\epsilon}.
\]



Las probabilidades proceden de las multiplicidades y de la
normalización de las emisiones APP--TRIT--TPK de
\cref{x_reciprocidad:iv:cont:medida,x_reciprocidad:nuclear:1}.
El registro identifica el estado precedente de cada emisión.
Si el cuadrado de transporte es
\(\mathcal J'_NR=(\bigoplus_aR_a)\mathcal J_N\), entonces



\begin{equation}
(\Lambda^pJ'_N)(\Lambda^pR)
=\Lambda^p(\bigoplus R_a)(\Lambda^pJ_N).
\label{x_reciprocidad:nuc11:eq:23}
\end{equation}



Se aplica
\(\Lambda^p(AB)=\Lambda^p(A)\Lambda^p(B)\).
Los cuadrados de las cinco operaciones correlativas y su límite
común son los de
\cref{x_reciprocidad:iv:rectangular-natural,x_reciprocidad:iv:cont:memoria-natural,x_reciprocidad:iv:cont:mapa-correlativo}; todas actúan en la misma construcción
del continuo.

En el caso positivo no estacionario,
\cref{x_reciprocidad:nuclear:10} demuestra
\(P_N^*P_N\downarrow A_\infty\) fuertemente y



\[
\mathcal Jv=(A_\infty^{1/2}v,(D_nP_nv)_{n\ge0}),
\qquad \mathcal J^*\mathcal J=I.
\]



Por densidad, \(\Lambda^p\mathcal J\) es isométrico para cada
grado finito \(p\).
La suma exterior se extiende a multiíndices de soporte finito
sobre la suma de Hilbert numerable; Parseval justifica el paso
al límite.
El terminal \(A_\infty^{1/2}\) y las memorias conservan
conjuntamente la forma inicial a profundidad arbitraria.

\subsection{Precisión de reconstrucción de áreas y volúmenes}

Para el lector \(M\), ordenemos los once valores propios
positivos de \(G\), con multiplicidad, como
\(\lambda_1\le\cdots\le\lambda_{11}\).
En los grados \(1\le p\le11\) de \(\mathbf1^\perp\),
la descomposición singular induce



\begin{equation}
\|(\Lambda^pM)^\dagger\|
=(\lambda_1\cdots\lambda_p)^{-1/2}.
\label{x_reciprocidad:nuc11:eq:24}
\end{equation}



Los modos exteriores utilizan \(p\) direcciones distintas.
El espectro calculado da las constantes óptimas
\(9/4\), \(81/16\), \(243\sqrt6/64\) y \(2187/128\)
para \(p=1,2,3,4\).
En grado doce, \(\Lambda^{12}M=0\); la recuperación del volumen
completo incorpora también la carga uniforme.
El producto tensorial de \(p\) lectores independientes sobre
media nula tiene cota \((9/4)^p\).
Tensorizar el lector por una identidad conserva la cota \(9/4\).
El registro completo \(\mathcal J\) tiene inversa de norma uno
sobre su imagen en cada grado.

\section{Memoria del transporte y holonomía relativa}

El acoplamiento reversible de \cref{x_reciprocidad:nuclear:10} es



\[
U(B,C)=\frac19
\begin{pmatrix}8B+C&\sqrt8(C-B)\\
\sqrt8(C-B)&B+8C\end{pmatrix}
=Q^*\operatorname{diag}(B,C)Q,
\]



con \(Q\) constante y ortogonal, fijada por la partición \(8+1\).
Por tanto,
\(U(B_2,C_2)U(B_1,C_1)=U(B_2B_1,C_2C_1)\),
conservando el orden de factores.
Sobre un lazo \(\gamma\), sean \(H_B,H_C\) los productos
ordenados de ambos transportes.
El bloque de memoria de la evolución completa es



\[
D_\gamma=\frac{\sqrt8}{9}(H_C-H_B).
\]



Se reconstruye exactamente



\begin{equation}
\boxed{H_B^{-1}H_C=I+\frac9{\sqrt8}H_B^{-1}D_\gamma.}
\label{x_reciprocidad:nuc11:eq:25}
\end{equation}



La igualdad es operatorial: registrar la respuesta del bloque
sobre una base recupera la holonomía relativa completa.
Su evaluación sobre un vector recupera la acción correspondiente.
Bajo cambio de carta en el punto base, ambos miembros se conjugan.
Si \(H_B=I\), el complemento determina \(H_C-I\).
Así se relacionan memoria y defecto de transporte sobre un
circuito; la realización diferencial utiliza adicionalmente
su conexión y el límite de circuitos.

\subsection{Un lazo elíptico–parabólico produce autoescala hiperbólica}

La construcción de \cref{x_reciprocidad:euler-trit-generadores} contiene el
Coxeter \(C\) y el nilpotente
\[
 N_0=\begin{pmatrix}0&1\\0&0\end{pmatrix}
\]
en sus portadores respectivos.
Se consideran sus realizaciones en un plano orientado común,
mediante las bases indicadas y la forma
\[
 \Omega=\begin{pmatrix}0&1\\-1&0\end{pmatrix}.
\]
El Coxeter \(C\) y el transporte \(I+N_0\) tienen determinante
uno y conservan \(\Omega\).
Esta carta permite componerlos conservando la distinción entre
los generadores de sus regímenes y sus métricas positivas propias.

Sea
\[
 P=I+N_0=\begin{pmatrix}1&1\\0&1\end{pmatrix}.
\]
Como \(N_0^2=0\), se cumple \(P^n=I+nN_0\) para todo
\(n\in\ZZ\).
El lazo ordenado de cuatro transportes da



\begin{equation}
H_n=C^{-1}P^{-n}CP^n
=\begin{pmatrix}1+n&n^2\\n&n^2-n+1\end{pmatrix}.
\label{x_reciprocidad:nuc11:eq:28}
\end{equation}



Para \(n\ne0\), su diferencia normalizada es



\[
F_n=\frac{H_n-I}{n}
=\begin{pmatrix}1&n\\1&n-1\end{pmatrix}.
\]



Multiplicando esas matrices se obtiene, sin introducir valores espectrales,



\begin{equation}
\boxed{F_n^2=H_n,\qquad F_n^2-nF_n-I=0.}
\label{x_reciprocidad:nuc11:eq:29}
\end{equation}



Con \(H_B=I\), la memoria es
\(D_{\gamma,n}=(\sqrt8/9)nF_n\).
En consecuencia,



\begin{equation}
\boxed{H_n=\left(\frac9{n\sqrt8}D_{\gamma,n}\right)^2.}
\label{x_reciprocidad:nuc11:eq:30}
\end{equation}



La respuesta matricial de memoria reconstruye la autoescala
hiperbólica. Para \(n>0\), el reconocimiento espectral posterior
identifica



\begin{equation}
\rho_n=\frac{n+\sqrt{n^2+4}}2,\quad
\operatorname{Spec}F_n=\{\rho_n,-\rho_n^{-1}\},\quad
\operatorname{Spec}H_n=\{\rho_n^2,\rho_n^{-2}\}.
\label{x_reciprocidad:nuc11:eq:31}
\end{equation}



El caso \(n=1\) produce el polinomio áureo.
La relación con el operador
\(F_{\rm av}=\begin{pmatrix}0&1\\1&1\end{pmatrix}\)
generado por el calendario en \eqref{x_reciprocidad:eq:fav} conserva el cambio
de carta:



\begin{equation}
F_1=PF_{\rm av}P^{-1},\qquad
H_1=PF_{\rm av}^{\,2}P^{-1}.
\label{x_reciprocidad:nuc11:eq:32}
\end{equation}



El transporte \(P\), de determinante uno, entrelaza ambas
matrices y conserva \(\Omega\).
El caso \(n=2\) produce las lecturas espectrales
\(1+\sqrt2\) y \(3+2\sqrt2\).
La composición relaciona estas realizaciones conservando sus
genealogías respectivas; la selección de una palabra de cuatro
factores permanece especificada como parte del lector.

En el mismo ejemplo \(n=1\), se obtiene el balance orientado



\begin{equation}
D_\gamma^{\mathsf T}\Omega D_\gamma=-\frac8{81}\Omega,
\qquad
T_\gamma^{\mathsf T}\Omega T_\gamma=\frac{89}{81}\Omega.
\label{x_reciprocidad:nuc11:eq:33}
\end{equation}



La suma de ambas contribuciones es \(\Omega\).
El término de memoria tiene orientación opuesta:
la contribución visible \(89/81\) se compensa con \(-8/81\).
Este balance alternante y el reparto positivo
\(19/27+8/27=1\) conservan sus formas respectivas.

La expresión de \(H_n\) vale para todo \(n\in\ZZ\).
La definición de \(F_n\) y la reconstrucción mediante su cuadrado
utilizan \(n\ne0\); la parametrización espectral anterior utiliza
\(n>0\).
En general,
\[
 D_{\gamma,n}^{\mathsf T}\Omega D_{\gamma,n}
 =-\frac{8n^2}{81}\Omega,\qquad
 T_{\gamma,n}^{\mathsf T}\Omega T_{\gamma,n}
 =\left(1+\frac{8n^2}{81}\right)\Omega.
\]
Las identidades \(N_0^2=0\) y \(C^2+C+I=0\) demuestran la
familia completa; el programa suplementario comprueba sus
especializaciones exactas.
La procedencia de esta formalización adicional se documenta
en el suplemento.

\section{Simetría, acción y carga conservada}

Sean \((E_n)_{n\in\ZZ}\) espacios vectoriales reales de dimensión
finita y \(U_n:E_n\to E_{n+1}\) los transportes invertibles de
la evolución completa.
La elevación cotangente actúa sobre
\(q_n\in E_n\) y \(p_n\in E_n^*\).
En bases duales, definimos el lagrangiano discreto auxiliar
\[
 L_n(q_n,q_{n+1},p_{n+1})
 =p_{n+1}^{\mathsf T}(q_{n+1}-U_nq_n).
\]
La acción se utiliza en sentido local:
para cada variación de soporte finito, se suma \(L_n\) sobre un
intervalo finito que contenga todos los términos afectados.
La primera variación es independiente de cualquier ampliación
de ese intervalo. En concreto,
\begin{equation}
 \delta\mathscr S
 =\sum_{n\in\ZZ}
 \left[
  \delta p_n^{\mathsf T}(q_n-U_{n-1}q_{n-1})
  +\delta q_n^{\mathsf T}(p_n-U_n^{\mathsf T}p_{n+1})
 \right],
 \label{x_reciprocidad:nuc11:accion-local}
\end{equation}
donde la suma es finita para cada variación admitida.
La estacionariedad respecto de todas estas variaciones equivale a
\[
 q_{n+1}=U_nq_n,\qquad
 p_n=U_n^{\mathsf T}p_{n+1}
 \quad(n\in\ZZ).
\]
La segunda ecuación es
\(p_{n+1}=U_n^{-\mathsf T}p_n\), es decir, la elevación
cotangente del transporte.

\begin{proposition}[Carga variacional del transporte entrelazado]
\label{x_reciprocidad:nuc11:noether}
Sea \(A_n\in\operatorname{End}(E_n)\) una familia que satisface
\[
 A_{n+1}U_n=U_nA_n.
\]
Sobre las soluciones estacionarias de
\eqref{x_reciprocidad:nuc11:accion-local}, la carga
\begin{equation}
 \boxed{J_n=p_n^{\mathsf T}A_nq_n=J_{n+1}}
 \label{x_reciprocidad:nuc11:carga}
\end{equation}
es constante.
\end{proposition}
\begin{proof}
Las ecuaciones de evolución y el entrelazamiento dan
\[
 J_{n+1}
 =p_{n+1}^{\mathsf T}A_{n+1}U_nq_n
 =p_{n+1}^{\mathsf T}U_nA_nq_n
 =p_n^{\mathsf T}A_nq_n.
\]
La conservación posee además una derivación variacional explícita.
Para parámetro constante \(\varepsilon\), las transformaciones
\[
 q_n\mapsto e^{\varepsilon A_n}q_n,\qquad
 p_n\mapsto e^{-\varepsilon A_n^{\mathsf T}}p_n
\]
conservan cada \(L_n\), porque
\(e^{\varepsilon A_{n+1}}U_n=U_ne^{\varepsilon A_n}\).
Si el parámetro se sustituye por una sucesión
\((\varepsilon_n)\) de soporte finito, la primera variación es
\[
 \delta L_n
 =(\varepsilon_{n+1}-\varepsilon_n)
   p_{n+1}^{\mathsf T}U_nA_nq_n.
\]
Sobre una solución, el último factor es \(J_n\).
La suma finita por partes da
\[
 \delta\mathscr S
 =\sum_{n\in\ZZ}\varepsilon_n(J_{n-1}-J_n).
\]
La estacionariedad para toda sucesión de soporte finito produce
\eqref{x_reciprocidad:nuc11:carga}.
\end{proof}

Para los lectores cíclicos anteriores se toma
\[
 A_0=S-S^{-1},\qquad
 \mathcal A=\operatorname{diag}(A_0,A_0)
\]
en el portador completo \(V\oplus V\).
La matriz \(\mathcal A\) conmuta con \(U(I,S^3)\) y
\(U(I,S^4)\), pues sus bloques son polinomios en \(S\).
En consecuencia, \(p_n^{\mathsf T}\mathcal A q_n\) conserva
una correlación orientada entre posiciones vecinas.

La construcción especifica un lagrangiano, su simetría y la carga
conservada correspondiente, en el sentido variacional desarrollado
en \cref{x_reciprocidad:nuclear:7}.
Su dominio es la realización cotangente de los lectores.
La acción física y sus unidades conservan los operadores y
normalizaciones de las construcciones correspondientes.
El lagrangiano auxiliar se obtiene después del transporte \(U_n\).
